% ----------------------------------------------------------------------
%  3.4.5. Диффузия и фазовые превращения
% ----------------------------------------------------------------------
\subsection{Диффузия и фазовые превращения в металлах и сплавах}\label{ssec:diff-phase}

\paragraph{Что лимитирует рост новой фазы}
Рост кристаллов новой фазы может ограничиваться двумя разными
процессами.
\begin{itemize}
\item \emph{Процессы на границе раздела} (перестройка решётки,
присоединение атомов). Если граница движется в матрице постоянного
среднего состава, скорость роста от времени не зависит~— \emph{линейный}
закон: $r\propto t$.
\item \emph{Диффузия} к растущему кристаллу (или от него). Тогда
положение границы меняется пропорционально $\sqrt t$~—
\emph{параболический} закон, как и длина диффузии $\sqrt{Dt}$.
\end{itemize}

\paragraph{Рост выделений при распаде твёрдого раствора}
Пусть из пересыщенного $\alpha$-раствора (концентрация $C^*$) растут
частицы $\beta$-фазы (концентрация $C_\beta$). Задачу упрощают:
число зародышей постоянно, они распределены равномерно, напряжения и
сопротивление границы не учитываются. Тогда весь материал можно
разбить на одинаковые сферические \emph{ячейки} радиуса $R$ (половина
среднего расстояния между зародышами) с растущей частицей радиуса
$r_0(t)$ в центре (рис.~\ref{fig:cell}).
\begin{shagi}
\item Между частицей и границей ячейки ($r_0<r<R$) концентрация
подчиняется уравнению диффузии в сферических координатах:
$\frac{\partial^2C}{\partial r^2}+\frac2r\frac{\partial C}{\partial r}=\frac1D\frac{\partial C}{\partial t}$.
\item У поверхности частицы~— равновесная растворимость: $C(r_0)=C_\alpha$.
\item Баланс вещества на движущейся границе:
$(C_\beta-C_\alpha)\frac{\d r_0}{\d t}=D\left.\frac{\partial C}{\partial r}\right|_{r_0}$.
\item На границе ячейки поток равен нулю (соседние частицы
«делят» растворённое вещество): $\partial C/\partial r|_R=0$.
\end{shagi}
В безразмерных переменных $\rho=r/R$, $\tau=Dt/R^2$ задача
универсальна и решается численно. Рост замедляется по мере того, как
раствор в ячейке обедняется до $C_\alpha$.

\begin{figure}[ht]
\centering
\begin{tikzpicture}
\begin{axis}[width=0.6\textwidth, height=5cm, xmin=0, xmax=1.05, ymin=0, ymax=1.15,
  xlabel={$r/R$}, ylabel={$C$}, xtick={1}, ytick={0.2,0.75,1.0}, yticklabels={$C_\alpha$,$C^*$,$C_\beta$},
  axis lines=left]
  \fill[fmAmber!30] (axis cs:0,0) rectangle (axis cs:0.25,1.0);
  \draw[fmAmber!70!black, thick] (axis cs:0,1.0) -- (axis cs:0.25,1.0) -- (axis cs:0.25,0.2);
  \addplot[fmBlue, thick, domain=0.25:1, samples=50] {0.2+0.53*(1-exp(-(x-0.25)/0.12))};
  \addplot[fmBlue!50, thick, dashed, domain=0.38:1, samples=50] {0.2+0.35*(1-exp(-(x-0.38)/0.2))};
  \draw[fmGray, dotted] (axis cs:0,0.75) -- (axis cs:1,0.75);
  \draw[fmGray, dotted] (axis cs:0,0.2) -- (axis cs:1,0.2);
  \node[font=\scriptsize, fmAmber!60!black] at (axis cs:0.125,0.6) {$\beta$};
  \node[font=\scriptsize, fmBlue, anchor=west] at (axis cs:0.6,0.67) {$t_1$};
  \node[font=\scriptsize, fmBlue!60, anchor=west] at (axis cs:0.75,0.45) {$t_2>t_1$};
\end{axis}
\end{tikzpicture}
\caption{Профиль концентрации в сферической ячейке вокруг растущей
частицы $\beta$-фазы (схема): у поверхности частицы~— равновесная
концентрация $C_\alpha$, вдали~— исходная $C^*$; со временем частица
растёт, а раствор обедняется}
\label{fig:cell}
\end{figure}

\paragraph{Реакционная диффузия: рост слоёв фаз}
Приведём в контакт два чистых металла, образующих промежуточные фазы. В
зоне диффузии вырастают \emph{слои} всех или части фаз, существующих на
изотермическом разрезе диаграммы состояния (рис.~\ref{fig:reactive}).
Важная особенность: по правилу фаз Гиббса в двойной системе при
постоянной температуре \emph{двухфазных областей} в зоне диффузии не
возникает. Поэтому профиль концентрации кусочно-непрерывен: внутри
каждого слоя концентрация меняется плавно, а на границах слоёв~—
скачком, причём граничные концентрации совпадают с границами фаз на
диаграмме, а величина скачка~— с шириной двухфазной области.

\begin{figure}[ht]
\centering
\begin{tikzpicture}
\begin{axis}[width=0.66\textwidth, height=5.2cm, xmin=-3, xmax=3, ymin=0, ymax=1.05,
  xlabel={$x$}, ylabel={$C_B$}, xtick=\empty, ytick={0,1}, yticklabels={A,B}, axis lines=left]
  \addplot[fmBlue, thick, domain=-3:-1.2, samples=30] {0.12*exp((x+1.2)/0.5)};
  \addplot[fmGreen!60!black, thick, domain=-1.2:0.1] {0.3+0.12*(x+1.2)/1.3};
  \addplot[fmRed, thick, domain=0.1:0.9] {0.58+0.08*(x-0.1)/0.8};
  \addplot[fmBlue, thick, domain=0.9:3, samples=30] {1-0.15*exp(-(x-0.9)/0.5)};
  \draw[fmGray, dashed] (axis cs:-1.2,0.12) -- (axis cs:-1.2,0.3);
  \draw[fmGray, dashed] (axis cs:0.1,0.42) -- (axis cs:0.1,0.58);
  \draw[fmGray, dashed] (axis cs:0.9,0.66) -- (axis cs:0.9,0.85);
  \node[font=\small] at (axis cs:-2.2,0.45) {$\alpha$};
  \node[font=\small] at (axis cs:-0.55,0.55) {$\beta$};
  \node[font=\small] at (axis cs:0.5,0.8) {$\gamma$};
  \node[font=\small] at (axis cs:2.0,0.6) {$\delta$};
  \draw[dimline] (axis cs:-1.2,0.95) -- node[above, font=\scriptsize]{$\Delta x_\beta$} (axis cs:0.1,0.95);
\end{axis}
\end{tikzpicture}
\caption{Профиль концентрации в зоне реакционной диффузии между
металлами A и B (схема): слои фаз $\alpha$, $\beta$, $\gamma$, $\delta$
разделены скачками концентрации, равными ширине двухфазных областей
диаграммы состояния}
\label{fig:reactive}
\end{figure}

Почему граничные концентрации остаются равновесными, хотя внутри слоёв
идёт диффузия? По модели Б.\,Я.~Пинеса, переход атома через межфазную
границу~— «внешняя» диффузия~— требует преодоления лишь
потенциального барьера, без образования вакансии, и поэтому идёт
гораздо быстрее «внутренней» диффузии в объёме фазы. Любое отклонение
от равновесия на границе мгновенно снимается, а скорость роста
лимитирует лишь диффузия внутри слоёв.

Толщина слоя каждой фазы растёт по \term{параболическому закону}
\begin{keyeq}\label{eq:parabolic}
(\Delta x_i)^2=K_i^2\,t ,
\end{keyeq}
который в большинстве случаев подтверждается опытом. Константа $K_i$
зависит не только от свойств самой фазы, но и от соседних фаз и от их
положения на диаграмме. Для расчёта используют средний коэффициент
взаимной диффузии в фазе
$\overline D_i=\frac1{\Delta C_i}\int_{C_i'}^{C_i''}\widetilde D\,\d C$.

Качественно скорость роста слоёв предсказывают \term{критерии Зайта}:
фаза растёт тем быстрее, чем больше её коэффициент взаимной диффузии и
чем меньше он у соседних фаз, и чем шире её область гомогенности и уже
соседние двухфазные области. Это важно на практике: интерметаллидные
прослойки сильно влияют на прочность соединений разнородных металлов.

\begin{podrobnee}[уравнение Кидсона]
Скорость смещения каждой границы определяется балансом потоков. Для
левой границы $i$-й фазы
$(C_i'-C_{i-1}'')\frac{\d x_i'}{\d t}=\widetilde D_i\left(\frac{\partial C}{\partial x}\right)_{C_i'}-\widetilde D_{i-1}\left(\frac{\partial C}{\partial x}\right)_{C_{i-1}''}$,
аналогично для правой; разность даёт уравнение Кидсона для толщины слоя
$\Delta x_i=x_i''-x_i'$. Подстановка Больцмана $\lambda=x/\sqrt t$ (градиенты
$\partial C/\partial x=t^{-1/2}\,\partial C/\partial\lambda$) приводит к
параболическому закону $\Delta x_i=2A\sqrt t$; по знаку $A$ судят, растёт
ли фаза ($A>0$) или не появляется вовсе ($A<0$). Главный вывод: скорость
роста слоя любой фазы связана со скоростями роста всех остальных фаз,
их диффузионными характеристиками, положением фаз на диаграмме и
составом исходной пары.
\end{podrobnee}

\paragraph{Распад пересыщенных твёрдых растворов}
Диффузия управляет всеми стадиями распада (старения).
\begin{itemize}
\item \emph{Начальная стадия}~— образование зон Гинье--Престона или
модулированной структуры при низких температурах. Её делают возможной
избыточные вакансии, закалённые с высокой температуры: комплексы
«вакансия~— растворённый атом» быстро мигрируют, а скопления вакансий
служат зародышами. Переход от быстрой стадии к медленной объясняют
уходом вакансий на стоки.
\item \emph{Конечная стадия}~— \emph{коагуляция} частиц стабильной фазы:
крупные частицы растут за счёт мелких, уменьшая суммарную поверхностную
энергию. Этот процесс тоже контролируется диффузией через раствор.
\end{itemize}

\paragraph{Спинодальный распад и «восходящая» диффузия}
Движущая сила~— градиент химического потенциала. Если свободная энергия
смешения $G^M(C)$ имеет участок с отрицательной кривизной,
$\partial^2G^M/\partial C^2<0$ (рис.~\ref{fig:spinodal},\,\textit{а}),
раствор там неустойчив к \emph{любым} малым флуктуациям состава и
непрерывно расслаивается на обогащённые и обеднённые области~— это
\term{спинодальный распад}.

Связь с диффузией видна из формулы Даркена. Для $G^M=RT(C_1\ln a_1+C_2\ln a_2)$
с учётом уравнения Гиббса--Дюгема ($C_1\,\d\ln f_1+C_2\,\d\ln f_2=0$)\ispr{В
оригинале (с.~483) уравнение Гиббса--Дюгема записано как
$C_1\ln f_1+C_2\ln f_2=0$; оно связывает дифференциалы:
$C_1\,\d\ln f_1+C_2\,\d\ln f_2=0$.}
\begin{equation}\label{eq:spinodal}
\frac{\partial^2G^M}{\partial C_i^2}=RT\left(\frac1{C_1}+\frac1{C_2}\right)
\left(1+\frac{\partial\ln f_i}{\partial\ln C_i}\right).
\end{equation}
\begin{glavnoe}
Внутри спинодали $\partial^2G^M/\partial C^2<0$, значит, отрицателен и
термодинамический множитель Даркена, а с ним и коэффициент взаимной
диффузии $\widetilde D$. Атомы текут \emph{против} градиента концентрации~—
\term{восходящая диффузия}, и флуктуации состава не рассасываются, а
растут. При обычном распаде (зарождение и рост) диффузия, наоборот,
нисходящая: зародыш растёт, «высасывая» растворённый элемент из
окружения (рис.~\ref{fig:spinodal},\,\textit{б, в}).
\end{glavnoe}

\begin{figure}[ht]
\centering
\begin{subfigure}[b]{0.32\textwidth}\centering
\begin{tikzpicture}
\begin{axis}[width=1.12\textwidth, height=4.6cm, xmin=0, xmax=1, ymin=-0.105, ymax=-0.02,
  xlabel={$C$}, ylabel={$G^M$}, xtick={0.26,0.74}, xticklabels={$C_{s1}$,$C_{s2}$}, ytick=\empty,
  axis lines=left, xticklabel style={font=\scriptsize}]
  \fill[fmRed!12] (axis cs:0.26,-0.105) rectangle (axis cs:0.74,-0.02);
  \addplot[fmBlue, thick, domain=0.005:0.995, samples=120] {x*ln(x)+(1-x)*ln(1-x)+2.6*x*(1-x)};
  \node[font=\scriptsize, fmRed] at (axis cs:0.5,-0.095) {спинодаль};
\end{axis}
\end{tikzpicture}
\caption{}
\end{subfigure}\hfill
\begin{subfigure}[b]{0.32\textwidth}\centering
\begin{tikzpicture}
\begin{axis}[width=1.12\textwidth, height=4.6cm, xmin=0, xmax=4, ymin=0, ymax=1,
  xlabel={$x$}, ylabel={$C$}, xtick=\empty, ytick=\empty, axis lines=left]
  \addplot[fmGray, dashed, domain=0:4] {0.5};
  \addplot[fmBlue!50, thick, domain=0:4, samples=120] {0.5+0.06*sin(deg(2*pi*x))};
  \addplot[fmBlue, thick, domain=0:4, samples=120] {0.5+0.2*sin(deg(2*pi*x))};
  \addplot[fmBlue!70!black, very thick, domain=0:4, samples=160] {0.5+0.35*tanh(4*sin(deg(2*pi*x)))};
\end{axis}
\end{tikzpicture}
\caption{}
\end{subfigure}\hfill
\begin{subfigure}[b]{0.32\textwidth}\centering
\begin{tikzpicture}
\begin{axis}[width=1.12\textwidth, height=4.6cm, xmin=0, xmax=4, ymin=0, ymax=1,
  xlabel={$x$}, ylabel={$C$}, xtick=\empty, ytick=\empty, axis lines=left]
  \addplot[fmGray, dashed, domain=0:4] {0.5};
  \addplot[fmRed, thick, samples=200, domain=0:4]
    {x<1.3 ? 0.5-0.25*exp(-(x-1.3)^2/0.25) : (x<1.7 ? 0.88 : (x<2.6 ? 0.5-0.25*exp(-(x-1.7)^2/0.25) : 0.5))};
  \node[font=\scriptsize, fmRed] at (axis cs:1.5,0.95) {частица};
\end{axis}
\end{tikzpicture}
\caption{}
\end{subfigure}
\caption{\textit{а}~— свободная энергия смешения раствора с тенденцией
к расслоению (расчёт для регулярного раствора); в области между
точками перегиба $C_{s1}$, $C_{s2}$ кривизна отрицательна;
\textit{б}~— спинодальный распад: синусоидальная модуляция состава
нарастает со временем (восходящая диффузия); \textit{в}~— зарождение
и рост: у частицы с резкой границей образуется обеднённая зона, диффузия
нисходящая}
\label{fig:spinodal}
\end{figure}

Распад, контролируемый восходящей диффузией, часто даёт правильную,
периодическую структуру чередующихся частиц.

С.\,Т.~Конобеевский показал, что восходящую диффузию могут вызывать и
\emph{напряжения}. Если атомные объёмы компонентов различны, то в
упруго изогнутом образце крупные атомы стремятся в растянутые слои,
мелкие~— в сжатые:
\begin{equation}\label{eq:konobeevsky}
\frac{\partial C}{\partial t}=D\frac{\partial^2C}{\partial x^2}-D_\varepsilon\frac{\partial^2\varepsilon}{\partial x^2},
\end{equation}
где $\varepsilon$~— упругая деформация, а $D_\varepsilon$ пропорционален
разнице атомных радиусов. Равновесие наступает, когда потоки
уравновешены. Напряжения, возникающие при выделении новой фазы,
ускоряют процесс выделения.

\paragraph{Диффузия в упорядочивающихся сплавах}
Упорядочение~— это диффузионное перераспределение атомов по узлам.
Равновесная степень дальнего порядка устанавливается в две стадии.
\begin{shagi}
\item В отдельных местах зарождаются и растут до соприкосновения
упорядоченные области с почти равновесным порядком. Так как в
неупорядоченном состоянии все узлы равноправны, «назначение» подрешёток
в разных областях может оказаться разным~— возникают \term{антифазные
домены}, и средний порядок по кристаллу ниже, чем внутри доменов.
\item Медленная стадия: одни домены растут за счёт других, пока порядок
не станет однородным.
\end{shagi}
Скорость зарождения упорядоченной фазы
$I=I_0\exp[-(Q_i+\Delta G_\text{кр}+\varepsilon)/kT]$ включает энергию
активации диффузии, работу образования критического зародыша и энергию
деформации матрицы. На графике Аррениуса для диффузии в
упорядочивающихся сплавах при температуре превращения $T_\text{пр}$
наблюдается излом, а в упорядоченной фазе зависимость нелинейна
(рис.~\ref{fig:order-arrh}): энергия активации растёт со степенью
дальнего порядка $\eta$, в простейшей модели
$Q_\text{уп}=Q\,(1+\text{const}\cdot\eta^2)$. В упорядоченных растворах
позиции атомов разных сортов неэквивалентны (подрешётки), и механизмы
диффузии очень сложны.

\begin{figure}[ht]
\centering
\begin{tikzpicture}
\begin{axis}[width=0.56\textwidth, height=4.8cm, xmin=0.8, xmax=1.6, ymin=-14, ymax=-7,
  xlabel={$1/T$}, ylabel={$\lg D^*$}, xtick={1.0}, xticklabels={$1/T_\text{пр}$}, ytick=\empty, axis lines=left]
  \addplot[fmBlue, thick, domain=0.8:1.0] {-8-5*(x-0.8)};
  \addplot[fmRed, thick, domain=1.0:1.6, samples=40] {-9-5*(x-1.0)-6*(x-1.0)^1.5};
  \addplot[fmBlue, dashed, domain=1.0:1.6] {-9-5*(x-1.0)};
  \draw[fmGray, dotted] (axis cs:1.0,-14) -- (axis cs:1.0,-7);
  \node[font=\scriptsize, fmBlue, anchor=south west] at (axis cs:0.82,-8.1) {беспорядок};
  \node[font=\scriptsize, fmRed, anchor=north east] at (axis cs:1.55,-12.6) {порядок};
\end{axis}
\end{tikzpicture}
\caption{Излом графика Аррениуса при упорядочении (схема): ниже
$T_\text{пр}$ энергия активации растёт вместе со степенью порядка}
\label{fig:order-arrh}
\end{figure}

\begin{itogi}
\begin{itemize}
\item Движущая сила диффузии~— градиент химического потенциала;
соотношения Онзагера $L_{ik}=L_{ki}$ связывают перекрёстные эффекты.
\item Законы Фика: $J=-D\,\partial C/\partial x$,
$\partial C/\partial t=D\,\partial^2C/\partial x^2$; длина диффузии
$\sim\sqrt{Dt}$; решения~— гауссиана (тонкий слой) и функция ошибок
(контакт двух тел); метод Больцмана--Матано даёт $\widetilde D(C)$.
\item Механизмы: вакансионный (главный для замещения), межузельный
(малые атомы), гантельный и краудионный (облучение).
$D=D_0e^{-Q/kT}$, $Q=E_f+E_m$; $\langle X^2\rangle=2Dt$; корреляционный
множитель $f\approx1-2/z$.
\item Пути ускоренной диффузии: поверхность $>$ границы $>$ дислокации
$>$ объём; модель Фишера разделяет граничный и объёмный потоки.
\item Эффект Киркендалла: $v=(D_1-D_2)\,\partial N_1/\partial x$,
$\widetilde D=N_2D_1+N_1D_2$; связь с самодиффузией через
термодинамический множитель.
\item Слои фаз растут параболически; при спинодальном распаде $\widetilde
D<0$~— диффузия восходящая.
\end{itemize}
\end{itogi}

\begin{proverka}
\item Почему коэффициент диффузии в законе Фика, вообще говоря, зависит
от концентрации, а подвижность~— гораздо слабее?
\item Во сколько раз нужно увеличить время цементации, чтобы толщина
слоя выросла в 3 раза?
\item Почему углерод диффундирует в железе на много порядков быстрее,
чем само железо?
\item Оцените корреляционный множитель для ОЦК-решётки и сравните с
точным значением.
\item По какому признаку на концентрационной кривой можно отличить
зернограничную диффузию от объёмной?
\item Почему метки в опыте Киркендалла смещаются в сторону латуни?
\item Почему при спинодальном распаде диффузия идёт «в гору»?
\end{proverka}
